\begin{frame}{일정과 마일스톤 — M1$\sim$M4가 제출물}
  \small
  \begin{tabular}{@{}p{0.52\textwidth}l@{}}
    \toprule
    \textbf{주 / 해야 할 일} & \textbf{마일스톤(제출)} \\
    \midrule
    \textbf{1}주: 환경 후보 2$\sim$3개 조사, \textbf{문제 정의 문서}(MDP
    5요소, 주 지표와 리턴의 이론적 범위), \textbf{베이스라인(무작위) 리턴
    측정} & \textbf{M1} 문제 정의 + 환경 선택 + 베이스라인 숫자 (1p) \\
    \textbf{2}주: Ch9$\sim$15 코드 재사용으로 알고리즘 $\geq$1개 구현,
    \textbf{시드 1개에서 엔드투엔드 동작 확인}(Pendulum 12만 스텝
    $\approx$20$\sim$30초), \textbf{학습 예산 추정} & \textbf{M2} 중간코드
    리뷰(코드 + 베이스라인 + e2e 숫자) \\
    \textbf{3}주: 나머지 시드 실행, 하이퍼파라미터 조정(모두 동일
    프로토콜), \textbf{학습 후 결정적 $a=\mu$ 평가}(시드마다 동일
    에피소드 수), 학습 곡선·시드 비교 표 도출 & \textbf{M3} 학습 곡선
    1장 + 시드 비교 표 \\
    \textbf{4}주: 수식적 정당화, 보고서·슬라이드 완성, \textbf{발표
    5$\sim$7분}, peer-review 2팀 & \textbf{M4} 보고서 + 발표 + 리뷰 2건 \\
    \bottomrule
  \end{tabular}
  \vskip0.8em
  \begin{block}{왜 M2에 ``학습 예산 추정''이 들어가는가}
    Block A(표 기반)에서는 1에피소드 1초라 예산이 문제가 안 됐지만, Block B
   에서는 Pendulum 12만 스텝이 CPU 22초인 반면 MuJoCo 대형과제의 100만 스텝은
    수 \textbf{시간} — \textbf{계산 전에 ``얼마나 걸리는가''를 계획하는 것
    자체가 방법론의 정직성}에 포함. (Pendulum 12만 스텝 $\approx$22초가
    ``정상 기준선''.)
  \end{block}
\end{frame}
